\documentclass[12pt]{amsart}
\usepackage[utf8]{inputenc}
\usepackage[T1]{fontenc}
\usepackage[french]{babel}
\usepackage{amsmath,amssymb,amsthm}
\usepackage{amscd}
\usepackage{mathrsfs}
\usepackage{array}
\usepackage{enumerate}

\theoremstyle{plain}
\newtheorem{prop}{Proposition}[section]
\newtheorem{thm}[prop]{Th\'eor\`eme}
\newtheorem{cor}[prop]{Corollaire}
\newtheorem{lem}[prop]{Lemme}

\theoremstyle{remark}
\newtheorem{rem}[prop]{Remarque}
\newtheorem{defn}[prop]{D\'efinition}

\DeclareMathOperator{\Hom}{Hom}
\DeclareMathOperator{\Ext}{Ext}
\DeclareMathOperator{\coker}{coker}
\DeclareMathOperator{\Spec}{Spec}
\DeclareMathOperator{\Pic}{Pic}
\DeclareMathOperator{\Alb}{Alb}
\DeclareMathOperator{\Gr}{Gr}
\DeclareMathOperator{\dimk}{dim}
\DeclareMathOperator{\lgB}{lg}
\DeclareMathOperator{\Tr}{Tr}
\DeclareMathOperator{\Ass}{Ass}

\newcommand{\cA}{\mathscr{A}}
\newcommand{\cB}{\mathscr{B}}
\newcommand{\cF}{\mathscr{F}}
\newcommand{\cG}{\mathscr{G}}
\newcommand{\cH}{\mathscr{H}}
\newcommand{\cO}{\mathscr{O}}
\newcommand{\cS}{\mathscr{S}}
\newcommand{\cE}{\mathscr{E}}
\newcommand{\cL}{\mathscr{L}}
\newcommand{\Ob}{\text{Ob}}
\newcommand{\Rf}{\mathbf{R}f_*}
\newcommand{\RHom}{\mathbf{R}\mathscr{H}\!om}
\newcommand{\Lf}{\mathbf{L}f^*}
\newcommand{\db}{D^b}
\newcommand{\dm}{D^-}
\newcommand{\dperfd}{D^b_{\text{parf}}}
\newcommand{\pr}{\text{pr}}

\title{Th\'eor\`eme de Lefschetz et crit\`eres de d\'eg\'en\'erescence de suites spectrales}
\author{P. Deligne}
\thanks{(1) Aspirant au F.N.R.S.}
\date{}

\begin{document}
\maketitle

\section{Th\'eor\`emes g\'en\'eraux}

Soient $\cA$ une cat\'egorie ab\'elienne et $\db(\cA)$ la sous-cat\'egorie pleine de sa cat\'egorie d\'eriv\'ee form\'ee des complexes born\'es. On rappelle que si $T$ est un foncteur cohomologique (voir \cite{Verdier2} ou \cite{Verdier1} p.~10) de $\db(\cA)$ dans une cat\'egorie ab\'elienne $\cB$, Verdier \cite{Verdier2} a d\'efini une suite spectrale, fonctorielle en $X \in \db(\cA)$:
\begin{equation}\label{eq:1.1}
E_2^{p,q} = T(H^q(X)[-p]) \Rightarrow T(X[p+q]).
\end{equation}

\begin{prop}\label{prop:1.2}
Soit $X \in \db(\cA)$; les conditions suivantes sont \'equivalentes:
\begin{enumerate}
\item[(i)] Quel que soit $T$ comme plus haut, la suite spectrale \eqref{eq:1.1} d\'eg\'en\`ere.
\item[(ii)] L'objet $X$ est isomorphe, dans $\db(\cA)$, \`a $\bigoplus_i H^i(X)[-i]$.
\end{enumerate}
\end{prop}

\begin{proof}
(i) $\Rightarrow$ (ii) D\'esignons par $T_i$ le foncteur cohomologique en $K$ d\'efini par
\[
T_i(K) = \Hom_{\db(\cA)}(H^i(X)[-i], K).
\]
Pour chacun des $T_i$ ($i \in \mathbb{Z}$), la suite spectrale \eqref{eq:1.1} s\'ecrit
\begin{equation}\label{eq:1.3}
E_2^{p,q} = \Ext_{\cA}^p(H^q(X), H^i(X)) \Rightarrow \Hom_{\db(\cA)}(H^i(X)[-i], K[p+q])
\end{equation}
et l\'homomorphisme de $\Hom_{\db(\cA)}(H^i(X)[-i], K[i])$ dans $E_2^{0,i} = \Hom_{\cA}(H^i(X), H^i(K))$ qui s\'en d\'eduit n\'est autre que celui provenant de la fonctorialit\'e de $H^0$.

Si les suites spectrales \eqref{eq:1.3} d\'eg\'en\`erent, les fl\`eches
\[
\Hom_{\db(\cA)}(H^i(X)[-i], K) \longrightarrow \Hom_{\cA}(H^i(X), H^i(K))
\]
sont surjectives. Par hypoth\`ese, tel est le cas si $K = X$, de sorte qu\'il existe des morphismes $a_i$ de $H^i(X)[-i]$ dans $X$, induisant l\'identit\'e sur le $i$-\`eme groupe de cohomologie. De plus, puisque $H^i(X) = 0$ sauf pour un nombre fini de valeurs de $i$, les $a_i$ sont nuls sauf un nombre fini, ce qui permet de d\'efinir
\[
a = \sum_i a_i : \bigoplus_i H^i(X)[-i] \longrightarrow X.
\]
Par construction, $a$ induit l\'identit\'e sur la cohomologie, donc est un isomorphisme.

(ii) $\Rightarrow$ (i) La suite spectrale \eqref{eq:1.1} est un foncteur additif en $X$, et d\'eg\'en\`ere lorsque $X$ n\'a qu\'un objet de cohomologie non nul, donc aussi si $X$ est somme de complexes n\'ayant qu\'un objet de cohomologie non nul.
\end{proof}

\begin{rem}\label{rem:1.4}
Il r\'esulte de la d\'emonstration pr\'ec\'edente qu\'il suffit de v\'erifier la condition (i) pour les foncteurs cohomologiques du type $\Hom_{\db(\cA)}(A, *)$ o\`u $A \in \Ob(\cA)$. Par dualit\'e, il suffirait aussi de consid\'erer les seuls foncteurs $\Hom_{\cA}(*, A)$, la condition (ii) \'etant autoduale.
\end{rem}

\begin{thm}\label{thm:1.5}
Soient $X \in \db(\cA)$, $n$ un entier et $u$ un endomorphisme de degr\'e $2$ de $X$,
\[
u : X \longrightarrow X[2],
\]
dont les it\'er\'es induisent des isomorphismes
\[
u^i : H^{n-i}(X) \xrightarrow{\sim} H^{n+i}(X) \qquad (i \geq 0).
\]
Sous ces conditions, $X$ est isomorphe \`a $\bigoplus_i H^i(X)[-i]$, dans la cat\'egorie $\db(\cA)$.
\end{thm}

\begin{proof}
V\'erifions la condition (i) de \eqref{prop:1.2}. Soit donc $T$ un foncteur cohomologique de $\db(\cA)$ dans $\cB$, et posons $T^p(K) = T(K[p])$.

On appelle partie primitive de $H^{n-i}(X)$, et on d\'esignera par ${}_0H^{n-i}(X)$, le noyau de l\'homomorphisme
\[
u^{i+1} : H^{n-i}(X) \longrightarrow H^{n+i+2}(X).
\]
On v\'erifie que pour $i \geq 0$, les applications
\begin{equation}\label{eq:1.6}
\sum_{k \geq 0} u^k : \bigoplus_{k \geq 0} {}_0H^{n-i-2k}(X) \longrightarrow H^{n-i}(X),
\qquad
\sum_{k \geq 0} u^k : \bigoplus_{k \geq 0} {}_0H^{n+i-2k}(X) \longrightarrow H^{n+i}(X)
\end{equation}
sont des isomorphismes.

La suite spectrale \eqref{eq:1.1} \'etant fonctorielle en $X$, et ``compatible'' aux translations sur les degr\'es, l\'endomorphisme $u$ de $X$ d\'efinit un endomorphisme de degr\'e $2$ de la suite spectrale
\[
u : {}\!'E_r^{p,q} \longrightarrow {}\!'E_r^{p,q+2}
\]
commutant aux diff\'erentielles $d_r$ ($r \geq 2$).

Lorsque $r = 2$ l\'endomorphisme $u$ de la suite spectrale
\[
u : T^p(H^q(X)) \longrightarrow T^p(H^{q+2}(X))
\]
se d\'eduit du morphisme $u : H^q(X) \to H^{q+2}(X)$ par fonctorialit\'e de $T^p$.

Soit $i \geq 0$. D\'esignons par ${}_0{}\!'E_r^{p,n-i}$ le noyau de l\'homomorphisme
\[
u^{i+1} : {}\!'E_r^{p,n-i} \longrightarrow {}\!'E_r^{p,n+i+2}
\]
(``partie primitive'' de la suite spectrale). Pour $r = 2$, les d\'ecompositions \eqref{eq:1.6} induisent par fonctorialit\'e de chaque $T^p$ des d\'ecompositions
\begin{equation}\label{eq:1.7}
\sum_{k \geq 0} u^k : \bigoplus_{k \geq 0} {}_0{}\!'E_r^{p,n-i-2k} \xrightarrow{\sim} {}\!'E_r^{p,n-i},
\qquad
\sum_{k \geq 0} u^k : \bigoplus_{k \geq 0} {}_0{}\!'E_r^{p,n+i-2k} \xrightarrow{\sim} {}\!'E_r^{p,n+i}.
\end{equation}

Prouvons par r\'ecurrence sur $r$ que les diff\'erentielles $d_r$ ($r \geq 2$) sont nulles. L\'hypoth\`ese de r\'ecurrence implique que $E_2 = E_r$, de sorte que la d\'ecomposition \eqref{eq:1.7} s\'applique au terme $E_r$ de la suite spectrale. Puisque $u$ commute aux $d_r$, il suffira de prouver que $d_r$ s\'annule sur la partie primitive de $E_r$.

Le diagramme suivant est commutatif:
\[
\begin{CD}
{}_0{}\!'E_r^{p,n-i} @>{d_r}>> {}\!'E_r^{p+r,n-i-r+1} \\
@V{\nu^{i+1}}VV @VV{\nu^{i+1}}V \\
{}\!'E_r^{p,n+i+2} @>{d_r}>> {}\!'E_r^{p+r,n-i+r+3}
\end{CD}
\]
La fl\`eche verticale gauche est nulle par d\'efinition de la primitivit\'e. La fl\`eche verticale droite est un monomorphisme, car la fl\`eche
\[
u^{i+r-1} : {}\!'E_r^{p+r,n-i-r+1} \longrightarrow {}\!'E_r^{p+r,n+i+r-1}
\]
est un isomorphisme et $r - 1 \geq 1$. Il en r\'esulte que $d_r$ est nul.
\end{proof}

\begin{rem}\label{rem:1.8}
On prendra garde qu\'on ne peut pas en g\'en\'eral choisir un isomorphisme entre $X$ et $\bigoplus_i H^i(X)[-i]$ qui transforme $u$ en la somme des applications $u : H^i(X) \to H^{i+2}(X)$. On obtient ais\'ement un contre-exemple en prenant $H^i(X) = 0$ pour $i \not= n$.

On peut d\'efinir un isomorphisme ``plus beau que les autres'' entre $X$ et $\bigoplus_i H^i(X)[-i]$, mais on n\'aura pas \`a s\'en servir ici.
\end{rem}

\begin{rem}\label{rem:1.9}
Soient $X \in \db(\cA)$, $n$ et $s$ des entiers et $u$ un endomorphisme de degr\'e $s$ (resp.~$2$) de $X$. On suppose que $H^i(X) = 0$ lorsque $i$ n\'est pas de la forme $ks$ ($0 \leq k \leq n$) (resp.~($0 \leq k \leq 2n$)) et que $u^{n-k}$ (resp.~$u^k$) induit un isomorphisme entre $H^{ks}(X)$ et $H^{(2n-k)s}(X)$ (resp.~entre $H^{n-k}(X)$ et $H^{n+k}(X)$). Sous ces hypoth\`eses, la d\'emonstration de \eqref{thm:1.5} montre encore que $X$ est isomorphe \`a $\bigoplus_i H^i(X)[-i]$.
\end{rem}

\begin{rem}\label{rem:1.10}
Soit $(X_i)_{i \in I}$ une famille d\'objets de $\db(\cA)$ et $u$ une famille d\'homomorphismes de degr\'e $2$: $u_i : X_i \to X_i[2]$. Si, quels que soient $i$ et $j$, $u_i$ induit un isomorphisme entre $H^{n-j}(X_i)$ et $H^{n+j}(X_i)$, il est encore vrai que chaque $X_i$ est isomorphe \`a la somme de ses objets de cohomologie; pour le voir, il suffit d\'appliquer formellement les raisonnements qui pr\'ec\`edent \`a
\[
u : \bigoplus_i X_i \longrightarrow \bigoplus_i X_i,
\]
m\^eme si cette somme n\'existe pas dans $\db(\cA)$.

La m\^eme remarque s\'applique \`a la variante \eqref{rem:1.9}.
\end{rem}

\begin{thm}\label{thm:1.11}
Soit $X \in \db(\cA)$ et supposons que $X$ admette des endomorphismes de degr\'e $0$, $\pi_i : X \to X$ tels que $H^j(\pi_i) = \delta_{ij}$. Alors, $X$ est isomorphe \`a $\bigoplus_i H^i(X)[-i]$ dans la cat\'egorie $\db(\cA)$.
\end{thm}

\begin{proof}
Appliquons encore le crit\`ere \eqref{prop:1.2}~(i). Le diagramme
\[
\begin{CD}
E_2^{p,q} @>{d_2}>> E_2^{p+2,q-1} \\
@V{\pi_q}VV @VV{\pi_{q-1}}V \\
E_2^{p,q} @>{d_2}>> E_2^{p+2,q-1}
\end{CD}
\]
est commutatif, la premi\`ere colonne est l\'identit\'e, et la seconde z\'ero, de sorte que $d_2 = 0$.
\end{proof}

\begin{cor}\label{cor:1.12}
Soit $X \in \db(\cA)$ la somme d\'une famille finie d\'objets $X_k \in \db(\cA)$. Si on a
\[
X \simeq \bigoplus_i H^i(X)[-i],
\]
alors, pour tout $k$, on a
\[
X_k \simeq \bigoplus_i H^i(X_k)[-i].
\]
\end{cor}

\begin{proof}
Soient $j_k$ l\'injection canonique de $X_k$ dans $X$, $\pr_k$ la projection canonique de $X$ sur $X_k$ et $\pi_i$ un endomorphisme de $X$ tel que $H^j(\pi_i) = \delta_{ij}$. Pour chaque $k$, les $\pr_k \pi_i j_k$ v\'erifient l\'hypoth\`ese de \eqref{thm:1.11}.
\end{proof}

\begin{rem}\label{rem:1.13}
On v\'erifie que pour que l\'isomorphisme entre $X$ et $\bigoplus_i H^i(X)[-i]$ puisse \^etre choisi tel que les $\pi_i$ s\'identifient aux $\pr_i$, il faut et il suffit que les $\pi_i$ soient des idempotents deux \`a deux orthogonaux. Il existe alors un seul isomorphisme ayant cette propri\'et\'e et induisant l\'identit\'e sur la cohomologie.
\end{rem}


\section{Applications}

Soit $f : X \to Y$ une application continue entre espaces topologiques (voire un morphisme de topos), $A$ un faisceau d\'anneaux sur $X$ et $\cF$ un faisceau de $A$-modules. Si $u \in H^2(X, A)$, $u$ d\'efinit un endomorphisme de degr\'e $2$ dans $\db(X)$, encore not\'e
\[
u : \cF \longrightarrow \cF[2].
\]
Par fonctorialit\'e, $u$ d\'efinit un endomorphisme de degr\'e $2$:
\[
Rf_*(u) : \Rf\cF \longrightarrow \Rf\cF[2].
\]
On dira que $(X, Y, u, \cF, n)$ v\'erifie la \emph{condition de Lefschetz}, ou que $\cF$ v\'erifie la condition de Lefschetz relativement \`a $u$, si, pour chaque $i \geq 0$, les fl\`eches obtenues en it\'erant $i$ fois la fl\`eche $Rf_*(u)$:
\[
Rf_*(u)^i : R^{n-i}f_*\cF \longrightarrow R^{n+i}f_*\cF \qquad (i \geq 0)
\]
sont des isomorphismes. Le th\'eor\`eme \eqref{thm:1.5} donne ici

\begin{prop}\label{prop:2.1}
Si $\cF$ v\'erifie la condition de Lefschetz relativement \`a $u$, alors, dans $\db(Y)$, on a
\[
\Rf\cF \simeq \bigoplus_i R^if_*\cF[-i]
\]
et en particulier, la suite spectrale de Leray
\[
H^p(Y, R^qf_*\cF) \Rightarrow H^{p+q}(X, \cF)
\]
d\'eg\'en\`ere.
\end{prop}

Si on y avait tenu, on aurait pu supposer que $\cF$ soit un $(A, f^*B)$-bimodule, o\`u $B$ est un faisceau d\'anneaux sur $Y$, et \'enoncer \eqref{prop:2.1} dans $\db(Y, B)$.

\subsection{}
Le cas des sch\'emas, munis de la topologie \'etale, ne rentre pas dans le cadre ci-dessus. On supposera pour la suite que la th\'eorie des $\mathbb{Z}_l$-faisceaux et $\mathbb{Q}_l$-faisceaux (non n\'ecessairement constructibles) ait \'et\'e faite, et que la cat\'egorie d\'eriv\'ee de la cat\'egorie de ces ``faisceaux'' soit raisonnable. En ce qui concerne les \'enonc\'es de d\'eg\'en\'erescence de suites spectrales, le lecteur qui n\'aurait pas confiance en le travail futur de Jouanolou pourra v\'erifier, en revenant \`a la d\'emonstration de (1.5), qu\'on peut s\'en passer.

Soit $l$ un nombre premier. On suppose pour toute la suite que $l$ est inversible sur tous les sch\'emas consid\'er\'es.

Rappelons que, sur tout sch\'ema $X$, le $\mathbb{Z}_l$-faisceau $\mathbb{Z}_l(1)$ est d\'efini par la formule
\[
\mathbb{Z}_l(1) = \varprojlim_n \boldsymbol{\mu}_{l^n}.
\]
Ce faisceau est un $\mathbb{Z}_l$-module inversible, ce qui permet de d\'efinir, pour tout $\mathbb{Z}_l$-faisceau ou $\mathbb{Q}_l$-faisceau $\cF$, le ``faisceau'' $\cF(i)$ par la formule
\[
\cF(i) = \mathbb{Z}_l(i) \otimes_{\mathbb{Z}_l} \cF \qquad (i \in \mathbb{Z}).
\]
On aura, pour tout morphisme $f : X \to Y$ de sch\'emas et tout couple de $\mathbb{Z}_l$ ou $\mathbb{Q}_l$-faisceaux $\cF$ et $\cG$ sur $X$ et $Y$ respectivement:
\begin{equation}\label{eq:2.3}
\Rf(\cF(i)) = (\Rf\cF)(i)
\end{equation}
et
\[
f^*\cG(i) = f^*\cG(i).
\]

Si $\cL$ est un faisceau inversible sur $X$, sa premi\`ere classe de Chern $l$-adique se trouve dans $H^2(X, \mathbb{Z}_l(1))$.

Si $f : X \to Y$ est un morphisme de sch\'emas, si $\cF$ est un $\mathbb{Z}_l$-faisceau (resp. $\mathbb{Q}_l$-faisceau) sur $X$ et si $u \in H^2(X, \mathbb{Z}_l(1))$ (resp. si $u \in H^2(X, \mathbb{Q}_l(1))$), on dira que $\cF$ v\'erifie la condition de Lefschetz relativement \`a $u$ si pour chaque $i \geq 0$, les fl\`eches obtenues en it\'erant $i$ fois la fl\`eche $Rf_*(u)$:
\[
Rf_*(u)^i : R^{n-i}f_*\cF \longrightarrow R^{n+i}f_*\cF(i) \qquad (i \geq 0)
\]
sont des isomorphismes. En vertu de (2.3), cela implique que toutes les fl\`eches
\[
Rf_*(u)^i : R^{n-i}f_*\cF(k) \longrightarrow R^{n+i}f_*\cF(k+i)
\]
sont des isomorphismes. Appliquant (1.10), on trouve:

\begin{prop}\label{prop:2.4}
Avec les notations pr\'ec\'edentes, si $\cF$ v\'erifie la condition de Lefschetz relativement \`a $u$, alors, dans $\db(Y, \mathbb{Z}_l)$ (resp. $\db(Y, \mathbb{Q}_l)$), on a
\begin{equation}\label{eq:2.5}
\Rf\cF \simeq \bigoplus_i R^if_*\cF[-i]
\end{equation}
et en particulier la suite spectrale de Leray
\[
H^p(Y, R^qf_*\cF) \Rightarrow H^{p+q}(X, \cF)
\]
d\'eg\'en\`ere.
\end{prop}

Si $g$ est un morphisme de sch\'emas de $Y$ dans $S$, (2.4) implique encore la d\'eg\'en\'erescence de la suite spectrale de Leray
\[
R^pg_*R^qf_*\cF \Rightarrow R^{p+q}(gf)_*\cF.
\]

\subsection{}
Reste \`a passer en revue quelques cas int\'eressants o\`u la condition de Lefschetz est v\'erifi\'ee; nous donnerons des d\'etails justificatifs dans (2.7). Soient donn\'es $f: X \to Y$, $u$, $\cF$ et $n$. Dans le cas topologique, on sait d\'apr\`es Godement \cite{Godement}, II, (4.7.1) que si $f$ est propre et $Y$ localement paracompact, la formation des images directes sup\'erieures commute au passage aux fibres, de sorte que la condition de Lefschetz se v\'erifie fibre par fibre. Dans le cadre des sch\'emas, la m\^eme r\'eduction s\'applique, la r\'ef\'erence \`a Godement \'etant remplac\'ee par une r\'ef\'erence \`a (SGA 4, XII, (5.1)). Dans le cadre des sch\'emas encore, si $f$ est propre et lisse et $\cF$ constant tordu ($=$ lisse, dans une nouvelle terminologie), les $R^if_*\cF$ sont encore constants tordus (cf. SGA 4, XVI, (2.2)) de sorte que si $Y$ est connexe, il suffit m\^eme de v\'erifier la condition de Lefschetz sur une fibre g\'eom\'etrique. Dans tous les exemples donn\'es, $f$ sera une application continue (resp. un morphisme de sch\'emas) propre et lisse de dimension relative topologique $2n$ (resp. alg\'ebrique $n$). Une application continue est dite lisse, si localement \`a la source elle est isomorphe \`a la projection de $\mathbb{R}^{2n} \times Y$ sur $Y$.

\begin{enumerate}
\item[(2.6.2)] Si $f$ est un morphisme propre et lisse de vari\'et\'es k\'ahl\'eriennes, le faisceau constant $\mathbb{C}$ v\'erifie la condition de Lefschetz relativement \`a la $2$-classe fondamentale de $X$.

\item[(2.6.3)] Si $X$ est un sch\'ema relatif (complexe) (voir \cite{Hakim}), projectif et lisse sur l\'espace topologique localement paracompact $Y$, et muni d\'un faisceau inversible relativement ample $\cO(1)$, le faisceau constant $\mathbb{C}$ v\'erifie la condition de Lefschetz relativement \`a la premi\`ere classe de Chern de $\cO(1)$.

\item[(2.6.4)] Supposons que $f$ soit un morphisme de sch\'emas projectif et lisse, avec $Y$ connexe et que soit donn\'e un faisceau inversible relativement ample $\cO(1)$ sur $X$. Si on peut remonter en caract\'eristique $0$ une fibre g\'eom\'etrique de $f$, et sa polarisation, alors le faisceau $\mathbb{Q}_l$-adique constant $\mathbb{Q}_l$ v\'erifie la condition de Lefschetz relativement \`a la premi\`ere classe de Chern $l$-adique de $\cO(1)$.

On sait qu\'on peut relever en caract\'eristique $0$ au sens pr\'ec\'edent dans les cas suivants:
\begin{itemize}
\item[$-$] $X$ est une vari\'et\'e de Severi-Brauer sur $Y$.
\item[$-$] $X$ est une intersection compl\`ete non singuli\`ere d\'hypersurfaces dans une vari\'et\'e de Severi-Brauer sur $Y$ (pour la polarisation induite par la polarisation canonique de la vari\'et\'e de Severi-Brauer).
\item[$-$] $X$ est un sch\'ema ab\'elien polaris\'e sur $Y$ (Mumford, \`a para\^itre).
\end{itemize}
On conjecture que (2.6.4) reste vrai lorsqu\'on ne sait pas se remonter en caract\'eristique $0$, mais ce n\'est prouv\'e que pour les surfaces.

\item[(2.6.5)] Si $f$ est un morphisme projectif et lisse de dimension relative $2$ et si $X$ est muni d\'un faisceau inversible relativement ample $\cO(1)$, alors, le faisceau constant $\mathbb{Q}_l$ v\'erifie la condition de Lefschetz relativement \`a la premi\`ere classe de Chern de $\cO(1)$.
\end{enumerate}

\subsection{}
Dans les cas (2.6.2) et (2.6.3), la $2$-classe de cohomologie donn\'ee sur $X$ induit sur les fibres de $f$ la $2$-classe d\'efinie par la structure k\'ahl\'erienne des fibres, de sorte que pour prouver (2.6.2) et (2.6.3), il suffit, d\'apr\`es (2.6.1) de prouver (2.6.2) dans le cas particulier o\`u $Y$ est r\'eduit \`a un point. L\'assertion r\'esulte alors de la th\'eorie de Hodge et de la d\'ecomposition de Hodge-Lepage (Weil \cite{Weil1}, IV, n${}^\circ$~6, cor. au Th.~5). La d\'emonstration donn\'ee dans \textit{loc. cit.} s\'applique encore si $\cF$ est un faisceau localement constant de $\mathbb{C}$-vectoriels de dimension finie, muni d\'une structure hermitienne d\'efinie positive localement constante.

D\'apr\`es (2.6.1), pour v\'erifier (2.6.4), il suffit de traiter le cas o\`u $Y$ est le spectre d\'un corps de caract\'eristique $0$, et le principe de Lefschetz permet de supposer $Y = \Spec(\mathbb{C})$.

On sait alors (cf. SGA 4, XI, (4.4)) que pour tout faisceau $\mathbb{Q}_l$-adique constant tordu $\cF$ sur $X$, si $\cF^{\text{an}}$ d\'esigne le faisceau localement constant de $\mathbb{Q}_l$-vectoriels sur l\'espace topologique $X^{\text{an}}$ d\'eduit de $\cF$, alors les $\mathbb{Q}_l$-vectoriels de cohomologie, alg\'ebrique et transcendante, de $\cF$ ou $\cF^{\text{an}}$, sont isomorphes
\begin{equation}\label{eq:2.7.2}
H^i(X, \cF) \simeq H^i(X^{\text{an}}, \cF^{\text{an}}).
\end{equation}
Choisissant un plongement de $\mathbb{Q}_l$ dans $\mathbb{C}$, on en d\'eduit des isomorphismes
\begin{equation}\label{eq:2.7.3}
H^i(X, \cF) \otimes_{\mathbb{Q}_l} \mathbb{C} \simeq H^i(X^{\text{an}}, \cF^{\text{an}} \otimes_{\mathbb{Q}_l} \mathbb{C}).
\end{equation}
Sur $\mathbb{C}$, l\'application exponentielle d\'efinit un isomorphisme dit canonique entre $\mathbb{Z}/l^n$ et $\boldsymbol{\mu}_{l^n}$, donc entre $\mathbb{Z}_l$ et $\mathbb{Z}_l(1)$. L\'image par l\'isomorphisme compos\'e
\[
H^2(X, \mathbb{Z}_l(1)) \longrightarrow H^2(X, \mathbb{Z}_l) \xrightarrow{\sim} H^2(X^{\text{an}}, \mathbb{Z}) \otimes \mathbb{Z}_l,
\]
de la premi\`ere classe de Chern $l$-adique de $\cO(1)$ co\"incide avec la premi\`ere classe de Chern transcendante de $\cO(1)$. L\'isomorphisme (2.7.3) ram\`ene donc le cas auquel on s\'etait r\'eduit de (2.6.4) \`a (2.6.3).

Pour la d\'emonstration de (2.6.5), voir Kleiman \cite{Kleiman}, pp.~28 et ss. Le point essentiel est le th\'eor\`eme suivant (Weil \cite{Weil2}, p.~127).

(2.8) Soit $S$ une surface projective lisse sur un corps $k$ et soit $D$ une courbe lisse trac\'ee sur $S$ telle que $\cO(D)$ soit ample. Alors, l\'application compos\'ee
\[
\Pic^0(X) \longrightarrow \Pic^0(D) = \Alb(D) \longrightarrow \Alb(X)
\]
est une isog\'enie.

\begin{rem}\label{rem:2.9}
Le th\'eor\`eme de d\'eg\'en\'erescence de suite spectrale (2.2) d\'eduit de (2.6.4) a \'et\'e conjectur\'e par Grothendieck, par des consid\'erations de ``poids'', lorsque $Y$ est projectif et lisse sur un corps alg\'ebriquement clos.
\end{rem}

\begin{rem}\label{rem:2.10}
Serre m\'a fait remarquer que les th\'eor\`emes de d\'eg\'en\'erescence de suites spectrales d\'eduits de (2.6.2) et (2.6.3) peuvent aussi se d\'emontrer par une extension facile de la m\'ethode utilis\'ee par Blanchard (\cite{Blanchard}, II, 1), tout au moins si on dispose de la dualit\'e de Poincar\'e sur la base. R\'eciproquement, la m\'ethode suivie ici permet de compl\'eter les th\'eor\`emes II (1.1) et II (1.2) de \textit{loc. cit.}, la d\'emonstration de l\'implication (ii) $\Rightarrow$ (i) \'etant essentiellement celle de Blanchard, qui devait supposer les faisceaux images directes sup\'erieures constants.
\end{rem}

\begin{thm}\label{thm:2.11}
Soient $f: X \to Y$ une application continue propre et lisse (2.6.1) d\'espaces topologiques de dimension relative $n$ avec $Y$ localement paracompact, $k$ un corps commutatif de caract\'eristique z\'ero et $u \in H^2(X, k)$. On suppose que
\begin{enumerate}
\item[(a)] Le faisceau $R^1f_*k$ est constant (i.e. simple).
\item[(b)] Les cup produits it\'er\'es
\[
u^i \wedge : H^0(Y, R^{n-i}f_*k) \longrightarrow H^0(Y, R^{n+i}f_*k)
\]
sont des isomorphismes.
\end{enumerate}
Alors, les conditions suivantes sont \'equivalentes:
\begin{enumerate}
\item[(i)] La classe $u$ est induite par un \'el\'ement de $H^2(X, k)$.
\item[(ii)] La transgression
\[
d_2 : H^0(Y, R^1f_*k) \longrightarrow H^2(Y, k)
\]
est nulle.
\item[(iii)] Dans $\db(Y, k)$, on a
\[
\Rf k \simeq \bigoplus_i R^if_*k[-i].
\]
\end{enumerate}
\end{thm}

L\'assertion (i) $\Rightarrow$ (iii) est contenue dans (2.1). L\'assertion (iii) implique la d\'eg\'en\'erescence de toute la suite spectrale de Leray
\begin{equation}\label{eq:2.12}
H^p(Y, R^qf_*k) \Rightarrow H^{p+q}(X, k),
\end{equation}
et a fortiori (ii), et il reste \`a prouver que (ii) $\Rightarrow$ (i).

Soit $x \in Y$ et $X_x = f^{-1}(x)$. La classe induite $u_x \in H^2(X_x, k)$ ne s\'annule sur aucune composante connexe de la vari\'et\'e $X_x$, qui est donc orientable. Si on d\'esigne par $\Tr_x$ l\'application compos\'ee
\[
H^n(X_x, k) \xrightarrow{u_x^{n/2}} H^0(X_x, k) = H_0(X_x, k) \longrightarrow k,
\]
la dualit\'e de Poincar\'e implique que la forme bilin\'eaire
\[
H^{n-i}(X_x, k) \times H^{n+i}(X_x, k) \longrightarrow k : (x, y) \mapsto \Tr_x(x \cup y)
\]
est une dualit\'e parfaite. Cette construction se ``globalise'' et fournit
\[
\Tr : R^nf_*k \xrightarrow{\sim} k.
\]

Prouvons tout d\'abord que dans (2.12), $d_2u \in H^2(Y, R^1f_*k)$ est $0$. Si $x \in H^0(Y, R^1f_*k)$, pour une raison de degr\'e, on a $x \cup u^n = 0$, donc
\[
0 = d_2(x \cup u^n) = d_2x \cup u^n - n x \cup u^{n-1} d_2u = -n x \cup u^{n-1} d_2u.
\]
D\'apr\`es (b), l\'application
\[
u^i \wedge : H^0(Y, R^1f_*k) \longrightarrow H^0(Y, R^{2n-1}f_*k)
\]
est un isomorphisme, de sorte que si (ii) est v\'erifi\'e, quel que soit $y \in H^0(Y, R^{2n-1}f_*k)$, on a
\[
y \cdot d_2u = 0,
\]
et en particulier
\begin{equation}\label{eq:2.13}
\Tr(y \cdot d_2u) = 0.
\end{equation}

Soit $V$ un $k$-vectoriel tel que $R^1f_*k$ soit un $k$-vectoriel isomorphe au faisceau constant $V$. On a $d_2u \in H^2(Y, V)$ et, d\'apr\`es (2.13), pour toute forme lin\'eaire $w$ sur $V$, l\'image par $w$ de $d_2u$ dans $H^2(Y, k)$ est nulle. Or on conclut que $d_2u = 0$.

Si $d_2u = 0$, la d\'emonstration de (1.5) montre que toutes les diff\'erentielles $d_r$ de (2.12) sont nulles, de sorte que $E_2 = E_\infty$.

Pour une raison de degr\'e, $d_2u^{n+1} = 0$, de sorte que $d_2u^{n+1} = (n+1) u^n \cup d_2u = 0$. Puisque $u^n \in H^0(Y, R^{2n}f_*k)$, on d\'eduit de (b) que $d_2u = 0$.

Pour une raison de degr\'e, on a $d_ru = 0$ pour $r > 3$, et $u$ provient donc d\'un \'el\'ement de $H^2(X, k)$.

(2.14) Dans le cadre des sch\'emas, le r\'esultat pr\'ec\'edent reste valable, mais n\'est utilisable que dans le cas ``g\'eom\'etrique'', plus pr\'ecis\'ement lorsque $\mathbb{Q}_l$ est isomorphe \`a $\mathbb{Q}_l(1)$ sur $Y$.

(2.15) Le corollaire (1.12) permet d\'etendre certains des \'enonc\'es qui pr\'ec\`edent \`a certains faisceaux non constants.

\begin{prop}\label{prop:2.16}
Soient $g$ et $f$ deux applications continues (resp. deux morphismes de sch\'emas) composables
\[
X \xrightarrow{f} Y \xrightarrow{g} Z,
\]
et $\cF$ un faisceau (resp. un $k$-faisceau ou $\mathbb{Q}_l$-faisceau) sur $X$. Supposons que
\begin{enumerate}
\item[(a)] $Rf_*\cF \in \db(Y)$ et $Rf_*\cF \simeq \bigoplus_i R^if_*\cF[-i]$.
\item[(b)] $R(gf)_*\cF \in \db(Z)$ et $R(gf)_*\cF \simeq \bigoplus_i R^i(gf)_*\cF[-i]$.
\end{enumerate}
Alors, pour tout $k$, on a
\[
R^kg_*(Rf_*\cF) \simeq \bigoplus_i R^kg_*(R^if_*\cF)[-i].
\]
\end{prop}

D\'apr\`es (b) le complexe
\[
R(gf)_*\cF \simeq Rg_*Rf_*\cF \simeq \bigoplus_k Rg_*(\bigoplus_i R^if_*\cF[-k])
\]
est somme de ses objets de cohomologie; d\'apr\`es (1.12), les $R^kg_*(Rf_*\cF)$, qui, \`a un d\'ecalage pr\'es, en sont facteurs directs gr\^ace \`a (a), jouissent de la m\^eme propri\'et\'e.

\begin{prop}\label{prop:2.17}
Soient $\cE$ un fibr\'e vectoriel localement libre sur un sch\'ema $S$, $\cF$ un faisceau $l$-adique constant tordu sur $\mathbf{P}(\cE)$ et $f$ la projection de $\mathbf{P}(\cE)$ sur $S$. Dans la cat\'egorie d\'eriv\'ee, on a
\[
\Rf\cF \simeq \bigoplus_i R^if_*\cF[-i].
\]
\end{prop}

L\'espace projectif est simplement connexe, et sa cohomologie est sans torsion (SGA 1, XI 1.1 et SGA 4, XVI 2.2). On a donc
\[
\cF = f^*f_*\cF
\]
et
\[
\Rf\cF = Rf_*\mathbb{Z}_l \otimes f_*\cF,
\]
formule qui permet de se ramener au cas o\`u $\cF = \mathbb{Z}_l$. Soit $u$ la classe de Chern du faisceau inversible $\cO(1)$ sur $\mathbf{P}(\cE)$. On sait alors que $\mathbb{Z}_l$ v\'erifie la condition de Lefschetz relativement \`a $u$, et on conclut par (2.2).

(2.18) De (2.17), on peut d\'eduire (\`a l\'aide de (2.16)) le m\^eme r\'esultat pour les fibr\'es en drapeaux de toutes esp\'eces d\'efinis par $\cE$. Dans le cadre topologique, ce qui pr\'ec\`ede s\'applique tel quel aux fibr\'es vectoriels sur $\mathbb{C}$, et s\'etend aux fibr\'es vectoriels sur $\mathbf{H}$ gr\^ace \`a (1.9). Pour les fibr\'es r\'eels, il faut se limiter au cas o\`u $\cF$ provient d\'un faisceau de $2$-torsion sur la base.

\begin{rem}\label{rem:2.19}
Soit $X$ un sch\'ema ab\'elien de dimension relative $g$ sur une base $Y$. Pour tout entier $n$, la multiplication par $n$, soit $[n]$, d\'efinit un endomorphisme
\[
[n]^* : \Rf\mathbb{Z}_l \longrightarrow \Rf\mathbb{Z}_l.
\]
Prenant des combinaisons lin\'eaires \`a coefficients rationnels (et \`a d\'enominateurs born\'es en terme de $g$) de ces endomorphismes, on construit des endomorphismes
\[
\pi_i : \Rf\mathbb{Z}_l \longrightarrow \Rf\mathbb{Z}_l
\]
tels que $H^j(\pi_i) = \delta_{ij}$. Appliquant (1.11), on conclut, sans hypoth\`ese projective, que
\[
\Rf\mathbb{Z}_l \simeq \bigoplus_i R^if_*\mathbb{Z}_l[-i].
\]
L\'id\'ee exploit\'ee ici est due \`a Lieberman.
\end{rem}

\section{La forme infinit\'esimale du th\'eor\`eme de semi-continuit\'e}

Je rappelle dans ce paragraphe quelques r\'esultats qu\'on parvient \`a trouver dans EGA, III, \S\,7. La formulation (3.4) m\'a \'et\'e signal\'ee par L.~Illusie.

(3.0) Dans tout ce paragraphe, on d\'esigne par $A$ un anneau local noeth\'erien fix\'e une fois pour toutes, par $m$ son id\'eal maximal et par $k$ son corps r\'esiduel. On d\'esigne par $\dm(A)$ la sous-cat\'egorie de la cat\'egorie d\'eriv\'ee de la cat\'egorie des $A$-modules, form\'ee des complexes born\'es sup\'erieurement de $A$-modules de type fini.

\begin{prop}\label{prop:3.1}
Soit $M$ un module de type fini sur $A$.

\textup{(i)} Pour tout homomorphisme (pas n\'ecessairement local) de $A$ dans un anneau artinien $B$, et tout $B$-module de type fini $N$, on a
\begin{equation}\label{eq:3.1.1}
\lgB_B(M \otimes N) \leq \lgB_B(N) \cdot \dim_k(M \otimes k).
\end{equation}

\textup{(ii)} Les conditions suivantes sur $M$ sont \'equivalentes:
\begin{enumerate}
\item[(a)] $M$ est un module libre.
\item[(b)] Quels que soient $B$ et $N$ comme en (i), on a
\begin{equation}\label{eq:3.1.2}
\lgB_B(M \otimes N) = \lgB_B(N) \cdot \dim_k(M \otimes k).
\end{equation}
\item[(c)] Quel que soit $n \in \mathbb{N}$, la condition (3.1.2) est v\'erifi\'ee pour $B = N = A/m^n$.
\item[(d)] Quel que soit $p \in \Ass(A)$, $M_p$ est plat sur $A_p$ et la condition (3.1.2) est v\'erifi\'ee pour $B = N = A_p/pA_p$.
\end{enumerate}
Si $A$ est sans composantes immerg\'ees, ces conditions \'equivalent encore \`a:
\begin{enumerate}
\item[(e)] Quel que soit le point maximal $p$ de $\Spec(A)$, on a
\begin{equation}\label{eq:3.1.2'}
\lgB_{A_p}(M_p) = \lgB(A_p) \cdot \dim_k(M \otimes k).
\end{equation}
\end{enumerate}
\end{prop}

La d\'emonstration (une application du lemme de Nakayama) est laiss\'ee au lecteur.

(3.2) Pour tout complexe $K$ et tout entier $n$, on d\'esignera par $\tau_{\leq n}(K)$ le complexe d\'efini par:
\begin{equation}\label{eq:3.2.1}
\tau_{\leq n}(K)^i = \begin{cases}
0 & \text{si } i < n \\
\coker(d^{n-1}) & \text{si } i = n \\
K^i & \text{si } i > n
\end{cases}
\end{equation}
de sorte que
\begin{equation}\label{eq:3.2.2}
H^i(\tau_{\leq n}(K)) = \begin{cases}
0 & \text{si } i < n \\
H^i(K) & \text{si } i \geq n.
\end{cases}
\end{equation}

Soit $K$ un objet de $\dm(A)$ (cf.~(3.0)). Quitte \`a remplacer $K$ par un complexe isomorphe dans la cat\'egorie d\'eriv\'ee, on peut supposer les composantes de $K$ libres de type fini.

Quels que soient $B$ et $N$ comme en (3.1) (i), on a, au niveau des complexes:
\[
\tau_{\leq n}(K \otimes N) \simeq \tau_{\leq n}(K) \otimes N.
\]
Prenant les caract\'eristiques d\'Euler-Poincar\'e des deux membres, on trouve
\[
\sum_{i \geq 0} (-1)^i \lgB_B(H^{n+i}(K \otimes N)) = \lgB_B(\coker(d^{n-1}) \otimes N) + \sum_{i > 0} (-1)^i \lgB_B(K^{n+i} \otimes N).
\]

Soustrayons de cette identit\'e l\'identit\'e analogue pour $B = N = k$, multipli\'ee par $\lgB_B(N)$. D\'apr\`es (3.1) (ii), a)$\Leftrightarrow$b), on a, les $K^i$ \'etant libres,
\begin{equation}\label{eq:3.2.3}
\sum_{i \geq 0} (-1)^i \lgB_B(H^{n+i}(K \otimes N)) - \lgB_B(N) \cdot \sum_{i \geq 0} (-1)^i \dim_k(H^{n+i}(K \otimes k))
\end{equation}
\[
= \lgB_B(\coker(d^{n-1}) \otimes N) - \lgB_B(N) \cdot \dim_k(\coker(d^{n-1}) \otimes k),
\]
ce qui permet d\'appliquer (3.1) au premier membre de (3.2.3).

\begin{thm}\label{thm:3.3}
\textup{(Th\'eor\`emes d\'echange et de semi-continuit\'e).} Soient $K \in \Ob \dm(A)$ \textup{(}cf.~(3.0)\textup{)} et $n \in \mathbb{Z}$.

\textup{(i)} Pour tout homomorphisme de $A$ dans un anneau artinien $B$ et tout $B$-module de type fini $N$, on a
\begin{equation}\label{eq:3.3.1}
\sum_{i \geq 0} (-1)^i \lgB_B(H^{n+i}(K \otimes N)) \leq \lgB_B(N) \cdot \sum_{i \geq 0} (-1)^i \dim_k(H^{n+i}(K \otimes k)).
\end{equation}

\textup{(ii)} Les conditions suivantes sur $K$ et $n$ sont \'equivalentes:
\begin{enumerate}
\item[(a)] Le foncteur cohomologique en le $A$-module $N$ : $H^n(K \otimes N)$ v\'erifie la propri\'et\'e d\'echange en $* = n$, i.e. les conditions \'equivalentes suivantes sur $K$ et $n$ sont v\'erifi\'ees:
\begin{enumerate}
\item[(a.1)] le foncteur $H^n(K \otimes N)$ est exact \`a gauche;
\item[(a.2)] il existe un module $Q_n$ tel que le foncteur (a.1) soit isomorphe au foncteur $\Hom_A(Q_n, N)$;
\item[(a.3)] le foncteur $H^{n-1}(K \otimes N)$ est exact \`a droite;
\item[(a.4)] il existe un module $P_n$ tel que le foncteur (a.3) soit isomorphe au foncteur $P_n \otimes N$;
\item[(a.5)] la fl\`eche $H^n(K) \to H^n(K \otimes k)$ est surjective;
\item[(a.6)] il existe un complexe born\'e sup\'erieurement $K'$, quasi-isomorphe \`a $K$ et \`a composantes plates (resp. libres de type fini), tel que $d^{n-1} = 0$.
\end{enumerate}
\item[(b)] Quels que soient $B$ et $N$ comme en (i), on a
\begin{equation}\label{eq:3.3.2}
\sum_{i \geq 0} (-1)^i \lgB_B(H^{n+i}(K \otimes N)) = \lgB_B(N) \cdot \sum_{i \geq 0} (-1)^i \dim_k(H^{n+i}(K \otimes k)).
\end{equation}
\end{enumerate}
Si $A$ est sans composantes immerg\'ees, ces conditions \'equivalent encore \`a:
\begin{enumerate}
\item[(c)] Quel que soit le point maximal $p$ de $\Spec(A)$, on a
\begin{equation}\label{eq:3.3.2'}
\sum_{i \geq 0} (-1)^i \lgB_{A_p}(H^{n+i}(K_p)) = \lgB(A_p) \cdot \sum_{i \geq 0} (-1)^i \dim_k(H^{n+i}(K \otimes k)).
\end{equation}
\end{enumerate}
\end{thm}

\begin{rem}\label{rem:3.3.3}
On e\^ut pu allonger la liste par les conditions analogues \`a (3.1) (ii) c) et d).
\end{rem}

L\'assertion (i) r\'esulte de (3.2.3) et (3.1) (i). D\'apr\`es EGA, III, (7.4.2), et sa d\'emonstration, les conditions (a.1), (a.3) et (a.6) sont \'equivalentes, et, si $K$ est repr\'esent\'e par un complexe born\'e sup\'erieurement de modules libres de type fini, elles signifient encore que $\coker(d^{n-1})$ est libre; compte tenu de (3.1) (ii) a)$\Leftrightarrow$b)$\Leftrightarrow$c) et de (3.2.3), ceci prouve l\'equivalence de (a.1), (a.3), (a.6), (b) et (c). On a, trivialement, (a.6)$\Rightarrow$(a.2)$\Rightarrow$(a.1) et (a.6)$\Rightarrow$(a.4)$\Rightarrow$(a.3). L'\'equivalence de la condition (a.5) avec les autres est contenue dans EGA, III, (7.5.2).

\begin{cor}\label{cor:3.4}
Soient $A$, $K$ et $n$ comme en (3.3), et $k \in \mathbb{N}$.

\textup{(i)} Quels que soient $B$ et $N$ comme en (3.3) (i), on a
\begin{equation}\label{eq:3.4.1}
\sum_{i=0}^{2k+1} (-1)^i \lgB_B(H^{n+i}(K \otimes N)) \leq \lgB_B(N) \cdot \sum_{i=0}^{2k+1} (-1)^i \dim_k(H^{n+i}(K \otimes k)).
\end{equation}

\textup{(ii)} Les conditions suivantes sur $K$, $n$ et $k$ sont \'equivalentes:
\begin{enumerate}
\item[(a)] Le foncteur cohomologique en $N$, $H^n(K \otimes N)$ v\'erifie la propri\'et\'e d\'echange en $* = n$ et en $* = n + 2k + 1$.
\item[(a$'$)] Le complexe $K$ est quasi-isomorphe \`a un complexe born\'e sup\'erieurement $K'$, \`a composantes plates (resp. libres de type fini) tel que $d^{n-1} = \ldots = d^{n+2k-1} = 0$.
\item[(b)] Quels que soient $B$ et $N$ comme en (i), on a \'egalit\'e dans (3.4.1).
\end{enumerate}
Si $A$ est sans composantes immerg\'ees, ces conditions \'equivalent encore \`a
\begin{enumerate}
\item[(c)] Quel que soit le point maximal $p$ de $\Spec(A)$, on a \'egalit\'e dans (3.4.1) pour $B = N = A_p$.
\end{enumerate}
\end{cor}

(3.4.2) Pour $k = 0$, il s\'agit l\`a de propri\'et\'es du seul hypertor $H^n(K \otimes N)$, consid\'er\'e comme foncteur en le $A$-module $N$. Elles s\'expriment encore en disant que $H^n(K)$ est libre et que, pour tout $N$, la fl\`eche canonique de $H^n(K) \otimes N$ dans $H^n(K \otimes N)$ est un isomorphisme.

(3.4.3) D\'autre part, si $K$ est un complexe parfait, on d\'eduit de (3.4) pour $n \ll 0$ un r\'esultat analogue \`a (3.3), la sommation sur $i \geq 0$ dans (3.3.1, 2, 2$'$) \'etant remplac\'ee par une sommation sur $i \leq 0$.

L\'assertion (i) et l\'equivalence de (a), (b) et (c) r\'esultent de (3.3) et de l\'identit\'e
\[
\sum_{i=0}^{2k+1} (-1)^i a_i = \sum_{i \geq 0} (-1)^i a_{n+i} + \sum_{i \geq 0} (-1)^i a_{n+2k+1+i}.
\]
Pour prouver (a$'$), on applique successivement (3.3) (ii) a)$\Leftrightarrow$a.6) \`a $(K, n)$ et \`a $(\tau_{\leq n+2k+1}(K), n+2k+1)$.

(3.5) Lorsque $A$ est artinien, (3.4) donne, pour $k = 0$, l\'in\'egalit\'e
\begin{equation}\label{eq:3.5.1}
\lgB_A H^n(K) \leq \lgB(A) \cdot \dim_k H^n(K \otimes k)
\end{equation}
et, si on a \'egalit\'e dans (3.5.1), alors la condition d\'echange est v\'erifi\'ee en $n$ et $n+1$, et $H^n(K)$ est un $A$-module libre.


\section{Morphisme de Gysin en cohomologie de Hodge}

Le paragraphe est r\'edig\'e d\'apr\`es des notes de A.~Grothendieck.

(4.0) Soit $f : X \to Y$ un morphisme propre entre sch\'emas $X$ et $Y$ lisses purement de dimension relative $n$ et $m$, sur une base $S$. On suppose que $S$ est noeth\'erien et admet un complexe dualisant pour pouvoir r\'ef\'erer \`a \cite{Hartshorne} pour la d\'efinition de $R f^!$; cette restriction est sans doute inutile. Enfin, posons $d = n - m$.

(4.1) Toute forme diff\'erentielle relative sur $Y$ d\'efinit par image r\'eciproque une forme diff\'erentielle relative sur $X$, d\'o\`u un morphisme
\begin{equation}\label{eq:4.1.1}
f^*\Omega_{Y/S}^\bullet = \Lf \Omega_{Y/S}^\bullet \longrightarrow \Omega_{X/S}^\bullet.
\end{equation}
D\'apr\`es la th\'eorie de la dualit\'e pour les morphismes lisses et la transitivit\'e de $R f^!$, on a canoniquement
\begin{equation}\label{eq:4.1.2}
R f^! \Omega_{Y/S}^\bullet = \Omega_{X/S}^\bullet[d].
\end{equation}
Appliquons \`a $\Omega_{X/S}^\bullet$ et $\Omega_{Y/S}^\bullet$ la formule d\'induction
\[
\Rf \mathbf{R}\mathscr{H}\!om(K, L) = \mathbf{R}\mathscr{H}\!om(\Lf K, R f^! L),
\]
qui s\'obtient ais\'ement \`a partir de la d\'efinition de $R f^!$ par bidualit\'e (Hartshorne \cite{Hartshorne}, chap.~VII, \S\,3). Puisque
\[
\Rf \mathbf{R}\mathscr{H}\!om(\Omega_{X/S}^\bullet, \Omega_{X/S}^\bullet) = \cO_Y,
\]
on trouve:
\[
\Rf \Omega_{X/S}^\bullet = \mathbf{R}\mathscr{H}\!om(\Lf \Omega_{Y/S}^\bullet, \Omega_{X/S}^\bullet)[d].
\]
La fl\`eche (4.1.1) d\'efinit donc par transposition une fl\`eche
\begin{equation}\label{eq:4.1.3}
u : \mathbf{R}\mathscr{H}\!om(\Omega_{X/S}^\bullet, \Omega_{X/S}^\bullet) \longrightarrow \Rf \Omega_{Y/S}^\bullet.
\end{equation}

D\'apr\`es la th\'eorie de la dualit\'e, se donner une fl\`eche (4.1.3) revient \`a se donner une fl\`eche
\begin{equation}\label{eq:4.1.4}
\Rf \Omega_{X/S}^\bullet \longrightarrow \Omega_{Y/S}^\bullet[-d].
\end{equation}

\begin{defn}\label{defn:4.2}
Soient $X$, $Y$, $f$, $S$, $n$, $m$ et $d$ comme en (4.0). On appelle morphisme de Gysin la fl\`eche (4.1.4)
\[
\Tr_f : \Rf \Omega_{X/S}^\bullet \longrightarrow \Omega_{Y/S}^\bullet[-d]
\]
et les fl\`eches telles que
\[
\Tr_f : H^p(X, \Omega_{X/S}^\bullet) \longrightarrow H^{p-2d}(Y, \Omega_{Y/S}^\bullet)
\]
qui s\'en d\'eduisent.
\end{defn}

\begin{prop}\label{prop:4.3}
Soit $f : X \to Y$ un morphisme propre et birationnel de sch\'emas lisses sur un corps $k$. Les fl\`eches
\[
f^* : H^p(Y, \Omega_{Y/k}^\bullet) \longrightarrow H^p(X, \Omega_{X/k}^\bullet)
\]
sont injectives.
\end{prop}

On se ram\`ene \`a supposer $X$ et $Y$ purement de dimension $n$. Le morphisme compos\'e
\[
\Tr_f \circ f^* : \Omega_{Y/k}^\bullet \longrightarrow \Rf \Omega_{X/k}^\bullet \longrightarrow R f_* f^* \Omega_{Y/k}^\bullet \longrightarrow \Omega_{Y/k}^\bullet
\]
est l\'identit\'e, car il co\"incide avec l\'identit\'e sur un ouvert dense. Appliquant le foncteur $H^p(Y, *)$, on trouve que le morphisme compos\'e
\[
\Tr_f \circ f^* : H^p(Y, \Omega_{Y/k}^\bullet) \longrightarrow H^p(X, \Omega_{X/k}^\bullet) \longrightarrow H^p(Y, \Omega_{Y/k}^\bullet)
\]
est l\'identit\'e, ce qui prouve (4.3).

\section{Cohomologie de Hodge et de De Rham}

(5.1) Soit $f : X \to S$ un morphisme lisse de sch\'emas. Par d\'efinition, les faisceaux de cohomologie de De Rham relative de $X$ sur $S$ sont donn\'es par
\[
\cH_{DR}^n(X/S) = R^n f_* \Omega_{X/S}^\bullet,
\]
o\`u $\Omega_{X/S}^\bullet$, le complexe de De Rham relatif, est un complexe diff\'erentiel $S$-lin\'eaire. On dispose d\'une suite spectrale
\begin{equation}\label{eq:5.1.1}
E_1^{p,q} = R^q f_* \Omega_{X/S}^p \Rightarrow \cH_{DR}^{p+q}(X/S).
\end{equation}

(5.2) Lorsque $S = \Spec(\mathbb{C})$ et que $X$ est propre et lisse, on sait par GAGA (\cite{Serre}) que la cohomologie de chacun des faisceaux coh\'erents $\Omega_{X/S}^p$, ainsi donc que l\'hypercohomologie de $\Omega_{X/S}^\bullet$ sont les m\^emes au sens alg\'ebrique (topologie de Zariski) ou transcendante (topologie usuelle de $X^{\text{an}}$), en particulier $\Omega_{X/S}^\bullet$ \'etant une r\'esolution du faisceau constant $\mathbb{C}$, on a
\[
\cH_{DR}^n(X) \simeq H^n(X^{\text{an}}, \mathbb{C})
\]
et la conjugaison complexe sur le faisceau constant $\mathbb{C}$ induit une conjugaison complexe (de nature transcendante) sur $H^n(X)$. La suite spectrale (5.1.1) s\'ecrit ici
\begin{equation}\label{eq:5.2.1}
E_1^{p,q} = H^q(X^{\text{an}}, \Omega^p) \Rightarrow H^{p+q}(X^{\text{an}}, \mathbb{C});
\end{equation}
on d\'esignera par $F^p(H^n(X))$ la filtration (d\'ecroissante) qu\'elle d\'efinit sur son aboutissement. D\'apr\`es Weil \cite{Weil1}, chap.~IV, n${}^\circ$~4, th.~2, la suite spectrale (5.2.1) peut se calculer comme suite spectrale du complexe double $H^0(X^{\text{an}}, \mathscr{A}^{\bullet\bullet})$ filtr\'e par le I er degr\'e, $\mathscr{A}^{p,q}$ d\'esignant le faisceau des formes diff\'erentielles $C^\infty$ bihomog\`enes de bidegr\'e $(p, q)$ (Weil \cite{Weil1}, chap.~II, n${}^\circ$~1).

Supposons maintenant que $X$ soit projective, donc $X^{\text{an}}$ une vari\'et\'e k\'ahl\'erienne. Les formes harmoniques forment alors un sous-complexe double (Weil \cite{Weil1}, chap.~II, n${}^\circ$~6, cor.~1 au th.~2) du complexe double pr\'ec\'edent, sur lequel $d'$ et $d''$ s\'annulent. L\'inclusion de ce sous-complexe induit un isomorphisme sur les termes $E_1$ des suites spectrales d\'efinies par ces doubles complexes, filtr\'es par le premier degr\'e, comme on le voit en consid\'erant l\'op\'erateur $H$ (composante harmonique), r\'etractation du complexe double $H^0(X^{\text{an}}, \mathscr{A}^{\bullet\bullet})$ sur le sous-complexe des formes harmoniques et v\'erifiant
\[
I - H = d'(2\delta' G) + (2\delta' G)d'
\]
(Weil \cite{Weil1}, chap.~IV, n${}^\circ$~1 : (I) et lemme~3; n${}^\circ$~3 : cor.~2 et chap.~II, n${}^\circ$~5, th.~2 (IX)).

Si on d\'esigne par $H^{p,q}(X)$ l\'image dans $H^{p+q}(X^{\text{an}}, \mathbb{C})$ de l\'espace des formes harmoniques de type $(p, q)$, on a donc
\begin{align}
H^q(X^{\text{an}}, \Omega^p) &\simeq H^{p,q}(X), \label{eq:5.2.2} \\
F^p(H^n(X)) &\simeq \bigoplus_{i \geq p} H^{i,n-i}(X), \label{eq:5.2.3} \\
H^n(X) &= \bigoplus_{p+q=n} H^{p,q}(X) \qquad \text{(conjugaison complexe)} \label{eq:5.2.4}
\end{align}
et la suite spectrale (5.2.1) d\'eg\'en\`ere.

Revenant au cas g\'en\'eral o\`u $X$ est seulement suppos\'e propre et lisse sur $\mathbb{C}$, on posera
\begin{equation}\label{eq:5.2.5}
H^{p,q}(X) = F^p(H^{p+q}(X)) \cap F^q(H^{p+q}(X)),
\end{equation}
de sorte que
\begin{equation}\label{eq:5.2.6}
H^{p,q}(X) \subset H^{p+q}(X).
\end{equation}
D\'apr\`es (5.2.3), cette notation est compatible avec celle utilis\'ee lorsque $X$ est projectif.

On posera encore
\begin{equation}\label{eq:5.2.7}
h^{p,q} = \dim H^q(X, \Omega_{X/S}^p).
\end{equation}

\begin{prop}\label{prop:5.3}
Soit $X$ un sch\'ema propre et lisse sur $\mathbb{C}$:
\begin{enumerate}
\item[(i)] La suite spectrale (5.2.1) d\'eg\'en\`ere,
\item[(ii)] On a
\begin{equation}\label{eq:5.3.1}
F^p(H^n(X)) \simeq \bigoplus_{i \geq p} H^{i,n-i}(X)
\end{equation}
et en particulier
\begin{equation}\label{eq:5.3.2}
H^n(X) = \bigoplus_{p+q=n} H^{p,q}(X) \qquad \text{(somme directe)}
\end{equation}
\begin{equation}\label{eq:5.3.3}
h^{p,q} = h^{q,p} = \dim H^{p,q}(X).
\end{equation}
\end{enumerate}
\end{prop}

\begin{proof}
On se ram\`ene \`a supposer $X$ connexe. Soit $N$ sa dimension.

D\'apr\`es le lemme de Chow et la r\'esolution des singularit\'es (Hironaka \cite{Hironaka}), il existe un morphisme projectif et birationnel $g : X' \to X$, avec $X'$ projectif et lisse sur $\mathbb{C}$. La suite spectrale
\begin{equation}\label{eq:5.3.4}
E_1^{p,q} = H^q(X, \Omega_{X}^p) \Rightarrow H^{p+q}(X)
\end{equation}
s\'envoie, par image r\'eciproque, dans la suite spectrale analogue
\begin{equation}\label{eq:5.3.5}
E_1^{p,q} = H^q(X', \Omega_{X'}^p) \Rightarrow H^{p+q}(X'),
\end{equation}
qui, d\'apr\`es (5.2) d\'eg\'en\`ere. D\'apr\`es (4.3), le morphisme $g^* : E_1 \to E_1$ est injectif. On en conclut, par r\'ecurrence sur $r \geq 1$, que $d_r = 0$, que $E_r = E_{r+1}$ et que $g^*$ injecte $E_r$ dans $E_r$.

Le diagramme
\[
\begin{CD}
E_r^{p,q} @>{d_r}>> E_r^{p+r,q-r+1} \\
@V{g^*}VV @VV{g^*}V \\
E_r^{p,q} @>{d_r}>> E_r^{p+r,q-r+1}
\end{CD}
\]
est commutatif. Puisque les suites spectrales (5.3.4) et (5.3.5) d\'eg\'en\`erent, et que $g^*$ est injectif sur leur terme initial, il l\'est encore sur leur aboutissement. Des formules (5.2.3), (5.2.4) appliqu\'ees \`a $X'$ r\'esulte que
\[
F^p(H^n(X')) \cap F^{n-p+1}(H^n(X')) = \{0\};
\]
on en d\'eduit la formule analogue
\begin{equation}\label{eq:5.3.6}
F^p(H^n(X)) \cap F^{n-p+1}(H^n(X)) = \{0\},
\end{equation}
car $g^*$ envoie $F^p(H^n(X))$ dans $F^p(H^n(X'))$ et commute \`a la conjugaison complexe. En particulier, on a
\[
\sum_{i \geq p} h^{i,n-i} + \sum_{i \geq n-p+1} h^{i,n-i} \leq \dim H^n(X),
\]
ce qu\'on peut \'ecrire
\begin{equation}\label{eq:5.3.7}
\sum_{i \geq p} h^{i,n-i} \leq \sum_{i \leq n-p} h^{i,n-i}.
\end{equation}
Par dualit\'e de Serre ($h^{i,j} = h^{N-i,N-j}$), l\'in\'egalit\'e (5.3.7) (pour $N = n$) implique l\'in\'egalit\'e oppos\'ee, de sorte que
\[
\dim F^p(H^n(X)) + \dim F^{n-p+1}(H^n(X)) = \dim H^n(X).
\]
D\'apr\`es (5.3.6), on a donc
\begin{equation}\label{eq:5.3.8}
H^n(X) = F^p(H^n(X)) \oplus F^{n-p+1}(H^n(X)).
\end{equation}
La d\'ecomposition (5.3.8) induit sur $F^{p-1}(H^n(X))$ (qui contient l\'un des facteurs) une d\'ecomposition
\[
F^{p-1}(H^n(X)) = F^p(H^n(X)) \oplus H^{p-1,n-p+1}(X),
\]
et la formule (5.3.1) s\'en d\'eduit par r\'ecurrence d\'ecroissante sur $p$. (5.3.2) et (5.3.3) r\'esultent aussit\^ot de (5.3.1) et (5.2.6), ce qui ach\`eve la d\'emonstration de (5.3).
\end{proof}

\begin{cor}\label{cor:5.4}
Sous les hypoth\`eses (5.3), on a:
\begin{enumerate}
\item[(i)] Une classe de cohomologie (complexe) $\alpha$ sur $X$ est dans $H^{p,q}(X)$ si et seulement si elle peut se repr\'esenter par une forme ferm\'ee $\alpha$ de type $(p, q)$.
\item[(ii)] Toute classe de cohomologie $\alpha$ peut se repr\'esenter par une forme $\alpha$ v\'erifiant $d'\alpha = d''\alpha = 0$.
\item[(iii)] Si la forme $\alpha$ v\'erifie $d'\alpha = d''\alpha = 0$, les conditions suivantes sont \'equivalentes:
\begin{enumerate}
\item[(a)] $(\exists \beta) \alpha = d'\beta$;
\item[(b)] $(\exists \beta) \alpha = d''\beta$;
\item[(c)] $(\exists \beta) \alpha = d\beta$.~(1)
\end{enumerate}
\end{enumerate}
\end{cor}

\footnotetext{(1) (Ajout\'e sur \'epreuves). On peut montrer que ces conditions \'equivalent encore \`a l\'existence d\'une forme $\beta$ telle que $\alpha = d'd''\beta$.}

Par d\'efinition, $\alpha \in H^{p,q}(X)$ si et seulement si existent dans la classe de $\alpha$ des formes ferm\'ees $\alpha_1$ et $\alpha_2$ dont les composantes de type $(p', q')$ sont nulles pour $p' \geq p$ (resp.~$q' \geq q$). Il existe alors $\beta$ tel que $\alpha_1 - \alpha_2 = d\beta$. Soit $\beta_1$ (resp.~$\beta_2$) la somme des composantes de $\beta$ de type $(p', q')$ pour $p' \geq p$ (resp.~$q' \geq q$); on a $\beta = \beta_1 + \beta_2$ et
\[
\alpha = \alpha_1 - d'\beta_1 = \alpha_2 + d''\beta_2
\]
est une forme ferm\'ee de type $(p, q)$ dans la classe de $\alpha$. Ceci prouve (i). Il suffit de prouver (ii) pour $\alpha$ bihomog\`ene, donc repr\'esent\'e par une forme bihomog\`ene ferm\'ee $\alpha$ d\'apr\`es (i). Une telle forme v\'erifie automatiquement $d'\alpha = d''\alpha = 0$.

Si $\alpha$ v\'erifie $d'\alpha = d''\alpha = 0$, ses composantes bihomog\`enes sont ferm\'ees, et chacune des conditions a), b), c) \'equivaut \`a la conjonction des conditions analogues pour les diverses composantes de $\alpha$, trivialement pour b) et c) et d\'apr\`es (5.3) (ii) et (5.4) (i) pour a). Supposons donc $\alpha$ ferm\'ee de type $(p, q)$ pour prouver (iii). La classe de cohomologie d\'efinie par $\alpha$ se trouve dans $H^{p,q}(X)$, donc dans $F^p(H^{p+q}(X))$ et est nulle si et seulement si elle se trouve d\'ej\`a dans $F^{p+1}(H^{p+q}(X))$, c\'est-\`a-dire si son image dans $H^q(X, \Omega^p)$ est nulle. Ceci prouve a)$\Leftrightarrow$c), et l\'equivalence a)$\Leftrightarrow$b) s\'en d\'eduit par conjugaison.

\begin{thm}\label{thm:5.5}
Soient $S$ un sch\'ema de caract\'eristique $0$ et $f : X \to S$ un morphisme propre et lisse. Alors
\begin{enumerate}
\item[(i)] Les faisceaux $R^q f_* \Omega_{X/S}^p$ sont localement libres de type fini de formation compatible \`a tout changement de base.
\item[(ii)] La suite spectrale (5.1.1)
\[
E_1^{p,q} = R^q f_* \Omega_{X/S}^p \Rightarrow \cH_{DR}^{p+q}(X/S)
\]
d\'eg\'en\`ere.
\item[(iii)] En chaque point de $S$, les faisceaux $R^q f_* \Omega_{X/S}^p$ et $R^p f_* \Omega_{X/S}^q$ ont m\^eme rang (sym\'etrie de Hodge).
\end{enumerate}
\end{thm}

La seconde assertion de (i) r\'esulte de la premi\`ere et de EGA, III, (7.8.5) appliqu\'e \`a chacun des $\Omega_{X/S}^p$; elle implique qu\'il suffit de v\'erifier (iii) pour $S$ spectre d\'un corps.

Des arguments standards permettent de se ramener successivement au cas o\`u $S$ est affine (car la question est locale sur $S$), noeth\'erien (par passage \`a la limite inductive d\'anneaux), spectre d\'un anneau local noeth\'erien (car la question est locale), spectre d\'un anneau local noeth\'erien complet (par fid\`ele platitude de la compl\'etion), spectre d\'un anneau local artinien. Pour effectuer cette derni\`ere r\'eduction, on repr\'esente l\'anneau local noeth\'erien complet $R$ comme limite projective de ses quotients artiniens $R/m^n$.

D\'apr\`es le th\'eor\`eme de comparaison EGA, III, (4.1.5), le terme $E_1$ de la suite spectrale (5.1.1) est limite projective des termes $E_1$ des suites spectrales analogues sur les $R/m^n$ de sorte que si celles-ci d\'eg\'en\`erent, (5.1.1) d\'eg\'en\`ere de m\^eme. Si l\'assertion (i) est vraie sur les $R/m^n$ ces termes $E_1$ forment un syst\`eme projectif de modules libres, se d\'eduisant les uns des autres par r\'eduction mod $m^n$ de sorte que leur limite projective est un module libre.

Supposons donc que $S = \Spec(A)$ avec $A$ anneau local artinien de corps r\'esiduel $k$ (de caract\'eristique $0$). Cet anneau admet un corps de repr\'esentants (EGA, O$_{IV}$, (19.8.6), (i) pour $W = k$, $u = \text{identit\'e}$) et peut donc \^etre vu comme une $k$-alg\'ebre locale de dimension finie. Le principe de Lefschetz permet de supposer $A = \mathbb{C}$, auquel cas (iii) r\'esulte de (i) et de (5.3.3).

Pour achever la d\'emonstration de (5.5), il reste donc \`a prouver la premi\`ere assertion de (i), et (ii), sous l\'hypoth\`ese suppl\'ementaire:
\begin{equation}\label{eq:5.5.1}
S = \Spec(A) \text{ pour une $\mathbb{C}$-alg\'ebre locale de dimension finie sur $\mathbb{C}$, $A$, d'id\'eal maximal $m$.}
\end{equation}
La suite spectrale (5.1.1) se r\'e\'ecrit
\begin{equation}\label{eq:5.5.2}
E_1^{p,q} = H^q(X, \Omega_{X/S}^p) \Rightarrow R^{p+q} f_* (\Omega_{X/S}^\bullet).
\end{equation}

\begin{lem}\label{lem:5.5.3}
Sous les hypoth\`eses (5.5) et (5.5.1), le complexe $\Omega_{X/S}^\bullet$ est une r\'esolution du faisceau constant $A$ sur l\'espace topologique $X^{\text{an}}$.
\end{lem}

Le complexe $\Omega_{X/S}^\bullet$ \'etant $A$-lin\'eaire, sa filtration $m$-adique est une filtration par des sous-complexes.

Pour cette filtration,
\[
\Gr_m \Omega_{X/S}^\bullet = \Gr_m(A) \otimes \Gr_0 \Omega_{X/S}^\bullet.
\]
Le complexe $\Gr_0 \Omega_{X/S}^\bullet$ est une r\'esolution du faisceau constant $\mathbb{C}$, car il n\'est autre que le complexe de De Rham de $X_k$. Le complexe $\Gr_m \Omega_{X/S}^\bullet$ est donc une r\'esolution de $\Gr_m(A)$, et l\'assertion en r\'esulte.

L\'espace $X^{\text{an}}$ est compact, de sorte que par GAGA (cf.~(5.2)) les hypercohomologies au sens alg\'ebrique et au sens transcendant de $\Omega_{X/S}^\bullet$ co\"incident, et, d\'apr\`es (5.5.3), co\"incident encore avec la cohomologie de $X^{\text{an}}$ \`a valeur dans $A$. D\`es lors,
\begin{equation}\label{eq:5.5.4}
R^n f_* (\Omega_{X/S}^\bullet) = H^n(X^{\text{an}}, A) = H^n(X^{\text{an}}, \mathbb{C}) \cdot \dim_\mathbb{C} A.
\end{equation}

D\'apr\`es (3.5.1) et EGA, III, (6.10.5), on sait que
\begin{equation}\label{eq:5.5.5}
\lgB_A H^q(X, \Omega_{X/S}^p) \leq \lgB(A) \cdot \dim_\mathbb{C} H^q(X_k^{\text{an}}, \Omega^p),
\end{equation}
l\'egalit\'e ne pouvant avoir lieu que si $H^q(X, \Omega_{X/S}^p)$ est $A$-libre.

On tire de (5.5.2) que
\begin{equation}\label{eq:5.5.6}
\sum_{p+q=n} \lgB_A H^q(X, \Omega_{X/S}^p) \geq \lgB_A R^n f_* (\Omega_{X/S}^\bullet),
\end{equation}
l\'egalit\'e ne pouvant avoir lieu pour tout $n$ que si la suite spectrale (5.5.2) d\'eg\'en\`ere.

Mettant bout \`a bout ces in\'egalit\'es, on trouve
\begin{equation}\label{eq:5.5.7}
\lgB(A) \cdot \sum_{p,q} \dim_\mathbb{C} H^q(X_k^{\text{an}}, \Omega^p) \geq \lgB(A) \cdot R^n f_* (\Omega_{X/S}^\bullet),
\end{equation}
l\'egalit\'e ne pouvant avoir lieu que si les conclusions (5.4) sont v\'erifi\'ees. Cette \'egalit\'e r\'esulte de (5.3) (i) appliqu\'e \`a $X_k^{\text{an}}$.

\section{Application \`a la cohomologie coh\'erente}

\begin{thm}\label{thm:6.1}
Soient $S$ un sch\'ema de caract\'eristique $0$ et $f : X \to S$ un morphisme projectif, lisse de dimension relative $n$. Quel que soit $p$ (par exemple $p = 0$), on a, dans $\db(S, \cO_S)$,
\[
\bigoplus_i R^i f_* \Omega_{X/S}^p[-i] \simeq \bigoplus_i R^i f_* \Omega_{X/S}^{p+1}[-i].
\]
\end{thm}

D\'esignons par $u$ la premi\`ere classe de Chern d\'un faisceau inversible relativement ample sur $X$, \`a valeur dans $H^1(X, \Omega_{X/S}^1)$. La classe $u$ d\'efinit par ``cup-produit'' des homomorphismes dans $\db(S)$:
\[
u \wedge : R^i f_* \Omega_{X/S}^p \longrightarrow R^{i+1} f_* \Omega_{X/S}^{p+1},
\]
et un endomorphisme de degr\'e $2$
\begin{equation}\label{eq:6.1.1}
u \wedge : \bigoplus_i R^i f_* \Omega_{X/S}^p[-i] \longrightarrow \bigl(\bigoplus_i R^i f_* \Omega_{X/S}^p[-i]\bigr)[2].
\end{equation}

\begin{lem}\label{lem:6.2}
L\'endomorphisme (6.1.1) v\'erifie les hypoth\`eses de (1.5).
\end{lem}

D\'apr\`es (5.4) (i) et le principe de Lefschetz, il suffit de d\'emontrer (6.2) lorsque $S = \Spec(\mathbb{C})$. L\'assertion r\'esulte alors du th\'eor\`eme de Lefschetz (voir (2.6.2) ou (2.6.3)) et de la d\'ecomposition de Hodge.

D\'apr\`es (1.5), le complexe $\bigoplus_i R^i f_* \Omega_{X/S}^p[-p]$ est isomorphe, dans la cat\'egorie d\'eriv\'ee, \`a la somme de ses objets de cohomologie, et il en va de m\^eme pour chacun des $R^i f_* \Omega_{X/S}^p$ qui, \`a un d\'ecalage pr\`es, en sont facteur direct (1.12).

\begin{thebibliography}{99}

\bibitem{Blanchard}
A.~Blanchard, \textit{Sur les vari\'et\'es analytiques complexes}, Ann.~Sc.~E.N.S., 73 (1956).

\bibitem{Godement}
R.~Godement, \textit{Th\'eorie des faisceaux}, Publ.~Inst.~math.~Univ.~de Strasbourg, XIII, Act.~Sci.~et~Ind., Hermann.

\bibitem{Hakim}
M.~Hakim, \textit{Sch\'emas relatifs}, th\`ese.

\bibitem{Hartshorne}
R.~Hartshorne, \textit{Residues and duality}, Lecture notes n${}^\circ$~20, Springer, 1966.

\bibitem{Hironaka}
H.~Hironaka, \textit{Resolution of singularities of an algebraic variety over a field of characteristic z\'ero}, Ann.~of Math., 79 (1964).

\bibitem{Kleiman}
S.~L.~Kleiman, \textit{Algebraic cycles and the Weil conjectures}, mimeographed notes, Columbia University.

\bibitem{Serre}
J.-P.~Serre, \textit{G\'eom\'etrie alg\'ebrique et g\'eom\'etrie analytique}, Ann.~Inst.~Fourier, Grenoble, 6 (1956), cit\'e GAGA.

\bibitem{Verdier2}
J.-L.~Verdier, \textit{Cat\'egories d\'eriv\'ees}, th\`ese \`a para\^itre \`a North Holl.~Publ.~Co.

\bibitem{Verdier1}
J.-L.~Verdier, \textit{Cat\'egories d\'eriv\'ees (\'etat 0)}, notes mim\'eographi\'ees, I.H.E.S.

\bibitem{Weil1}
A.~Weil, \textit{Introduction \`a l'\'etude des vari\'et\'es k\'ahl\'eriennes}, Publ.~Inst.~math.~Univ.~de Nancago, VI, Act.~Sci.~et~Ind., Hermann.

\bibitem{Weil2}
A.~Weil, \textit{Sur les crit\`eres d'\'equivalence en g\'eom\'etrie alg\'ebrique}, Math.~Ann., 128 (1954), p.~95--127.

\bibitem{SGA4}
\textbf{SGA 4}, M.~Artin, A.~Grothendieck et J.-L.~Verdier, \textit{Cohomologie \'etale des sch\'emas}. S\'eminaire de g\'eom\'etrie alg\'ebrique de l'I.H.E.S. (1963--1964), \`a para\^itre \`a North Holl.~Publ.~Co.

\bibitem{SGA5}
\textbf{SGA 5}, A.~Grothendieck, \textit{Cohomologie $l$-adique et rationalit\'e des fonctions $L$}, S\'eminaire de g\'eom\'etrie alg\'ebrique de l'I.H.E.S. (1964--1965).

\end{thebibliography}

\vspace{1cm}
\noindent \textit{Manuscrit re\c{c}u le 1\textsuperscript{er} juillet 1968.}

\end{document}
